\chapter{Robust H-infinity Control}

\section{The H-infinity Norm of LTI Systems}
In practical engineering systems, mathematical models are never perfect. Parameter variations, unmodeled high-frequency dynamics, and physical wear introduce unstructured perturbations. The $H_\infty$ norm is the key mathematical tool in robust control theory used to quantify a system's worst-case performance and guarantee robust stability in the presence of these model uncertainties.

Consider the stable, continuous-time LTI system:
\begin{align}
    \dot{x}(t) &= A x(t) + B_w w(t) \label{eq:hinf_state} \\
    z(t) &= C_z x(t) + D_{zw} w(t) , \label{eq:hinf_output}
\end{align}
where $w(t) \in \R^m$ represents worst-case external disturbances, and $z(t) \in \R^r$ is the controlled output performance vector.

\subsection{Frequency-Domain and Time-Domain Definitions}
In the frequency domain, the $H_\infty$ norm of the transfer function $G(s) = C_z (sI - A)^{-1} B_w + D_{zw}$ is defined as the peak gain (or maximum singular value) of the system over all frequencies:
\begin{equation}
    \| G \|_\infty = \sup_{\omega \in \R} \sigma_{\max}\left( G(j\omega) \right) , \label{eq:hinf_freq}
\end{equation}
where $\sigma_{\max}(\cdot)$ denotes the maximum singular value of the matrix.

\begin{figure}[H]
    \centering
    \includegraphics[width=9.5cm]{hinf_norm.png}
    \caption{H-infinity norm defined as the peak singular value over all frequencies.}
    \label{fig:hinf_norm}
\end{figure}

In the time domain, the $H_\infty$ norm corresponds to the \textbf{induced $L_2$-to-$L_2$ norm} of the system. Let $L_2$ be the space of square-integrable signals:
\begin{equation}
\| w \|_{L_2} = \left( \int_0^\infty \| w(t) \|^2 dt \right)^{1/2} < \infty .
\end{equation}
Assuming zero initial conditions ($x(0) = 0$), the $H_\infty$ norm is defined as:
\begin{equation}
    \| G \|_\infty = \sup_{w \neq 0, w \in L_2} \frac{\| z \|_{L_2}}{\| w \|_{L_2}} = \sup_{w \neq 0, w \in L_2} \frac{\left( \int_0^\infty \| z(t) \|^2 dt \right)^{1/2}}{\left( \int_0^\infty \| w(t) \|^2 dt \right)^{1/2}} . \label{eq:hinf_time}
\end{equation}
This indicates the worst-case disturbance amplification ratio. For any disturbance signal $w(t)$, the output energy is guaranteed to satisfy:
\begin{equation}
\| z \|_{L_2} \le \| G \|_\infty \| w \|_{L_2} .
\end{equation}

\section{The Bounded Real Lemma (BRL)}
Evaluating the $H_\infty$ norm using the frequency-domain definition \eqref{eq:hinf_freq} requires sweeping across all frequencies, which is computationally cumbersome. The \textbf{Bounded Real Lemma (BRL)} provides a state-space characterization of the $H_\infty$ norm, formulating the inequality $\| G \|_\infty < \gamma$ as a single Linear Matrix Inequality.

\begin{mytheorem}{Continuous-Time Bounded Real Lemma}
The LTI system \eqref{eq:hinf_state}--\eqref{eq:hinf_output} is asymptotically stable and satisfies $\| G \|_\infty < \gamma$ if and only if there exists a symmetric positive definite matrix $P \in \R^{n \times n}$ satisfying:
\begin{equation}
    \begin{bmatrix} A^\trp P + P A & P B_w & C_z^\trp \\ B_w^\trp P & -\gamma I & D_{zw}^\trp \\ C_z & D_{zw} & -\gamma I \end{bmatrix} \prec 0 . \label{eq:brl_lmi}
\end{equation}
\end{mytheorem}

\begin{proof}
We prove the sufficiency ($\Leftarrow$) of the BRL. Assume there exists $P \succ 0$ satisfying the LMI \eqref{eq:brl_lmi}. We define a quadratic Lyapunov (or storage) function:
\begin{equation}
V(x(t)) = x(t)^\trp P x(t) \ge 0 .
\end{equation}
Differentiating $V(x(t))$ along the system trajectories $\dot{x} = A x + B_w w$ yields:
\begin{align}
    \dot{V}(x) &= \dot{x}^\trp P x + x^\trp P \dot{x} \nonumber \\
    &= (A x + B_w w)^\trp P x + x^\trp P (A x + B_w w) \nonumber \\
    &= x^\trp (A^\trp P + P A) x + 2 x^\trp P B_w w . \label{eq:brl_vdot}
\end{align}
To establish $\| G \|_\infty < \gamma$, we seek to show that for any non-zero disturbance $w(t) \in L_2$ under $x(0) = 0$, we have:
\begin{equation}
\| z \|_{L_2}^2 < \gamma^2 \| w \|_{L_2}^2 \implies \int_0^\infty \left( \| z(t) \|^2 - \gamma^2 \| w(t) \|^2 \right) dt < 0 .
\end{equation}
Let us define the Hamiltonian-like dissipation term:
\begin{equation}
    \mathcal{D}(x, w) = \dot{V}(x) + \frac{1}{\gamma} \| z(t) \|^2 - \gamma \| w(t) \|^2 . \label{eq:brl_dissipation}
\end{equation}
Substitute $z(t) = C_z x(t) + D_{zw} w(t)$ and $\dot{V}(x)$ from \eqref{eq:brl_vdot} into \eqref{eq:brl_dissipation}:
\begin{align}
    \mathcal{D}(x, w) &= x^\trp (A^\trp P + P A) x + 2 x^\trp P B_w w \nonumber \\
&\quad + \frac{1}{\gamma} (C_z x + D_{zw} w)^\trp (C_z x + D_{zw} w) - \gamma w^\trp w .
\end{align}
We can write this quadratic form compactly as:
\begin{equation}
\mathcal{D}(x, w) = \begin{bmatrix} x \\ w \end{bmatrix}^\trp \begin{bmatrix} A^\trp P + P A + \frac{1}{\gamma} C_z^\trp C_z & P B_w + \frac{1}{\gamma} C_z^\trp D_{zw} \\ B_w^\trp P + \frac{1}{\gamma} D_{zw}^\trp C_z & \frac{1}{\gamma} D_{zw}^\trp D_{zw} - \gamma I \end{bmatrix} \begin{bmatrix} x \\ w \end{bmatrix} .
\end{equation}
For the system to satisfy the dissipation inequality $\mathcal{D}(x, w) < 0$ for all $\begin{bmatrix} x^\trp & w^\trp \end{bmatrix}^\trp \neq 0$, the central matrix must be strictly negative definite:
\begin{equation}
    \begin{bmatrix} A^\trp P + P A + \frac{1}{\gamma} C_z^\trp C_z & P B_w + \frac{1}{\gamma} C_z^\trp D_{zw} \\ B_w^\trp P + \frac{1}{\gamma} D_{zw}^\trp C_z & \frac{1}{\gamma} D_{zw}^\trp D_{zw} - \gamma I \end{bmatrix} \prec 0 . \label{eq:brl_proof_schur_pre}
\end{equation}
We can rewrite this matrix by separating the terms involving $C_z$ and $D_{zw}$:
\begin{equation}
\begin{bmatrix} A^\trp P + P A & P B_w \\ B_w^\trp P & -\gamma I \end{bmatrix} + \frac{1}{\gamma} \begin{bmatrix} C_z^\trp \\ D_{zw}^\trp \end{bmatrix} \begin{bmatrix} C_z & D_{zw} \end{bmatrix} \prec 0 .
\end{equation}
Applying the Schur complement to this expression immediately yields the BRL LMI:
\begin{equation}
\begin{bmatrix} A^\trp P + P A & P B_w & C_z^\trp \\ B_w^\trp P & -\gamma I & D_{zw}^\trp \\ C_z & D_{zw} & -\gamma I \end{bmatrix} \prec 0 .
\end{equation}
Now, integrating $\mathcal{D}(x, w) < 0$ from $t = 0$ to $t = \infty$:
\begin{equation}
\int_0^\infty \dot{V}(x(t)) dt + \frac{1}{\gamma} \int_0^\infty \| z(t) \|^2 dt - \gamma \int_0^\infty \| w(t) \|^2 dt < 0 .
\end{equation}
Since $x(0) = 0$ and the system is stable ($\lim_{t\to\infty} x(t) = 0$), the first term evaluates to:
\begin{equation}
\int_0^\infty \dot{V}(x(t)) dt = V(x(\infty)) - V(x(0)) = 0 .
\end{equation}
Thus, we have:
\begin{equation}
\frac{1}{\gamma} \| z \|_{L_2}^2 - \gamma \| w \|_{L_2}^2 < 0 \implies \| z \|_{L_2}^2 < \gamma^2 \| w \|_{L_2}^2 .
\end{equation}
This completes the proof of sufficiency.
\end{proof}

\subsection{Discrete-Time Bounded Real Lemma}
For discrete-time systems $x_{k+1} = A x_k + B w_k$, $z_k = C_z x_k + D_{zw} w_k$, the discrete-time BRL characterizes the performance $\| G \|_\infty < \gamma$ using the following LMI.

\begin{mytheorem}{Discrete-Time Bounded Real Lemma}
The discrete LTI system is asymptotically stable and satisfies $\| G \|_\infty < \gamma$ if and only if there exists a symmetric positive definite matrix $P \in \R^{n \times n}$ satisfying:
\begin{equation}
    \begin{bmatrix} A^\trp P A - P & A^\trp P B_w & C_z^\trp \\ B_w^\trp P A & B_w^\trp P B_w - \gamma I & D_{zw}^\trp \\ C_z & D_{zw} & -\gamma I \end{bmatrix} \prec 0 . \label{eq:brl_dt_lmi}
\end{equation}
\end{mytheorem}

\begin{figure}[H]
    \centering
    \includegraphics[width=9.5cm]{hinf_norm_discrete.png}
    \caption{Discrete-time H-infinity norm maximum singular value frequency plot.}
    \label{fig:hinf_norm_discrete}
\end{figure}

\section{Robust H-infinity State Feedback Synthesis}
We now design a state feedback gain $K \in \R^{m \times n}$ for the system:
\begin{align}
    \dot{x}(t) &= A x(t) + B_u u(t) + B_w w(t) \\
z(t) &= C_z x(t) + D_{zu} u(t) ,
\end{align}
such that the closed-loop system under the control law $u(t) = -K x(t)$ is stable and satisfies $\| G_{cl} \|_\infty < \gamma$. We assume $D_{zw} = 0$ to focus on the essential LMI structure.

The closed-loop dynamics are:
\begin{align}
    \dot{x}(t) &= (A - B_u K) x(t) + B_w w(t) \\
z(t) &= (C_z - D_{zu} K) x(t) .
\end{align}
Applying the continuous-time BRL \eqref{eq:brl_lmi} with $D_{zw} = 0$, the closed loop satisfies $\| G_{cl} \|_\infty < \gamma$ if there exists $P \succ 0$ satisfying:
\begin{equation}
    \begin{bmatrix} (A - B_u K)^\trp P + P (A - B_u K) & P B_w & (C_z - D_{zu} K)^\trp \\ B_w^\trp P & -\gamma I & 0 \\ C_z - D_{zu} K & 0 & -\gamma I \end{bmatrix} \prec 0 . \label{eq:hinf_bmi}
\end{equation}
This is a BMI due to the products of the decision variables $P$ and $K$.

To linearize the inequality, we perform a congruence transformation:
\begin{enumerate}
    \item Define $W = P^{-1} \succ 0$.
    \item Pre- and post-multiply \eqref{eq:hinf_bmi} by the block diagonal matrix $\text{diag}(W, I, I)$:
    \begin{equation}
\begin{bmatrix} W & 0 & 0 \\ 0 & I & 0 \\ 0 & 0 & I \end{bmatrix} \begin{bmatrix} (A - B_u K)^\trp P + P (A - B_u K) & P B_w & (C_z - D_{zu} K)^\trp \\ B_w^\trp P & -\gamma I & 0 \\ C_z - D_{zu} K & 0 & -\gamma I \end{bmatrix} \begin{bmatrix} W & 0 & 0 \\ 0 & I & 0 \\ 0 & 0 & I \end{bmatrix} \prec 0 .
    \end{equation}
    This simplifies to:
    \begin{equation}
\begin{bmatrix} W (A - B_u K)^\trp + (A - B_u K) W & B_w & W (C_z - D_{zu} K)^\trp \\ B_w^\trp & -\gamma I & 0 \\ (C_z - D_{zu} K) W & 0 & -\gamma I \end{bmatrix} \prec 0 .
    \end{equation}
    \item Define the auxiliary variable $Y = K W$, which yields:
    \begin{equation}
\begin{bmatrix} A W + W A^\trp - B_u Y - Y^\trp B_u^\trp & B_w & W C_z^\trp - Y^\trp D_{zu}^\trp \\ B_w^\trp & -\gamma I & 0 \\ C_z W - D_{zu} Y & 0 & -\gamma I \end{bmatrix} \prec 0 .
    \end{equation}
\end{enumerate}

This leads directly to the robust $H_\infty$ synthesis theorem.

\begin{mytheorem}{H-infinity State Feedback LMI Synthesis}
The robust $H_\infty$ state feedback gain $K$ that minimizes the worst-case disturbance amplification $\gamma$ is solved by the following convex semidefinite program:
\begin{equation}
\min_{W, Y} \gamma \quad \text{subject to} .
\end{equation}
\begin{align}
    W &\succ 0 \\
    \begin{bmatrix} A W + W A^\trp - B_u Y - Y^\trp B_u^\trp & B_w & W C_z^\trp - Y^\trp D_{zu}^\trp \\ B_w^\trp & -\gamma I & 0 \\ C_z W - D_{zu} Y & 0 & -\gamma I \end{bmatrix} &\prec 0 . \label{eq:hinf_sf_lmi}
\end{align}
The optimal feedback gain is reconstructed as:
\begin{equation}
K^* = Y^* (W^*)^{-1} .
\end{equation}
\end{mytheorem}
If the system matrices $[A \quad B_u]$ are subject to polytopic uncertainty, we can solve the LMI \eqref{eq:hinf_sf_lmi} simultaneously at all vertices using the single common variables $W$ and $Y$. This guarantees robust stability and a guaranteed $H_\infty$ attenuation bound $\gamma$ for all plants in the uncertainty polytope.
